\chapter*{Conclusion}
\addcontentsline{toc}{chapter}{Conclusion}
\markboth{CONCLUSION}{CONCLUSION}
\clearpage

%si ambroise critique le povuoir personnelle, l'ambition etc, (voir chapitre 1) sa relation avec les empereurs montre plutto qu'il est au coeur de ce système d'ambition et d'accaperement de l'autorité
%

Au terme d'une étude approfondie des homélies, exégèses, traités, et correspondances d'Ambroise de Milan, il apparait clairement que l'évêque ne s'est jamais seulement placé en théoricien observant le pouvoir impérial à distance, sans pour autant n'être qu'un acteur pragmatique de son temps. Au contraire, il concilie parfaitement les deux dimensions. C'est ce qui marque sans aucun doute le plus l'ensemble de l'analyse : la capacité d'Ambroise à lier sa vie épiscopale et politique avec son exercice théorique. Toute l'originalité de sa pensée, comme nous avons tenté de le démontrer, réside dans sa capacité à élaborer un nouveau modèle du pouvoir et du gouvernement romain à partir des fondements théoriques et réalistes qui composent sa vie épiscopale. Il s'appuie sur l'ensemble de l'héritage romain, pour y ancrer la doctrine chrétienne afin de définir l'essence et la légitimité du pouvoir impérial.

\bigskip

L'objectif de cette recherche était justemment de comprendre avec précision le procédé intelectuel et politique d'Ambroise, l'ayant amené à réfléchir, avec plus ou moins de réussite, sur l'exercice du pouvoir. Plusieurs fondements de sa culture sont clairement ressortis. D'un côté l'héritage de l'aristocrate et magistrat romain, éduqué à la philosophie classique et aux souvenirs de la République romaine. De l'autre, un évêque intransigeant, convaincu que les Écritures peuvent apporter une conduite et une morale idéales à la société. Face à un pouvoir monarchique fort, incarné par des figures d'autorité importante comme Théodose, ou au contraire plus jeunes et instables comme Valentinien~II, l'ambition d'Ambroise n'était pas de s'attaquer à la toute puissance de l'\latin{imperium}, mais de proposer un modèle de justice qui subordonne moralement et inconsciemment le pouvoir impérial. Toute sa réflexion consiste à questionner ce que peut être la société humaine, et donc chrétienne, idéale afin de la projeter sur le monde romain qu'il côtoie pour le faire évoluer. Il remet ainsi en question les différentes formes d'autorité pour s'en approprier, il prend du recul sur la légitimité de l'héritage dynastique, et impose aux empereurs l'exigence de la moralité chrétienne qu'il relie à la défense de la foi et de l'Église. Tout cet aspect s'inscrit également dans une volonté de réussite de la carrière personnelle pour Ambroise. Il se permet ces questionnements et théories car il est au coeur des institutions impériales et qu'il intervient auprès des empereurs comme conseiller. Ses connaissances des codes aristocratiques et politiques lui permettent de se détacher parmi les évêques et de jouer ce double rôle religieux et politique. Parfois figure de contestation, parfois de soutien de la légitimité impériale, il apparaît dans tous les cas comme une figure incontournable du monde politique romain.

\bigskip

De plus, sa pensée politique l'amène forcément à transformer l'autorité et le rôle de l'Église dont il fait partie. Ses théories et actions n'impactent pas uniquement sa carrière ni celles des empereurs mais l'ensemble du corps épiscopal. Il essaie de donner à l'Église leur prore sphère d'action où l'\latin{auctoritas} de l'empereur ne peut pas dépasser l'autorité des évêques. Les questions dogmatiques et de l'administration religieuse sont uniquement du ressors de l'Église. Ambroise demande ainsi aux empereurs un rôle uniquement de défense de la «~vraie foi~» à travers la sécurité des évêques et la restriction des cultes polythéistes et hérétiques. Les travaux de pensée politique d'Ambroise de Milan ont donc amené un renouveau dans la relation entre l'Église et l'Empire. L'Église chrétienne a continué son intégration aux structures administratives de l'Empire tout en se construisant comme une autorité morale autonome, derrière la figure ambrosienne.

\bigskip

La réponse à cette problématique sur la place de l'autorité d'Ambroise et son influence sur les divers empereurs romains réside dans le concept d'encadrement moral. En effet, face à un régime monarchique qui semble difficile à dépasser, Ambroise fait le choix de l'accepter pour mieux le conditionner et le faire évoluer vers sa vision du gouvernement idéal. Conscient que la stabilité de l'Empire a facilité la mission universaliste du christianisme, Ambroise adopte un certain réalisme politique pour centrer sa réfléxion sur l'exercice du pouvoir par les empereurs. Il cherche ainsi à imposer son idéologie : la justice dans l'Empire n'est plus seulement décidé par la volonté seule de l'empereur ou par la force des armes, c'est la morale chrétienne, la piété et le rejet de la coercition qui guide le pouvoir terrestre.

\bigskip

Ainsi, entre théorie utopique et actions réalistes, Ambroise met en place sa vision chrétienne du pouvoir et de la société. Son influence ne lui permet pas encore d'imposer ses volontés sur chaque action politique qui implique une affaire religieuse, mais il est certain qu'il incarne une volonté chrétienne de la diplomatie et de la force du pouvoir romain.

\bigskip

Le premier grand thème qui a traversé cette recherche est celui de la nécessaire conciliation entre un idéal politique à l'ambition collective et la réalité pragmatique d'un pouvoir personnel fort qui est difficile à dépasser. C'est ce qui marque le plus clairement la théologie politique ambrosienne : le compromis toujours trouvé entre idéal et réalité atteignable. Les premières descriptions théoriques vont dans ce sens. La conception par l'évêque de Milan du pouvoir romain s'est aussi bien fait par l'observation des sociétés animales que par une analyse approfondie de l'exercice du pouvoir. Il peut ainsi calquer différement modèles sur sa projection du gouvernement idéal et tirer le meilleur de chaque étude.

\bigskip

Chaque texte d'Ambroise évoquant une question politique présente un objectif commun : l'obtention d'un consensus général, et plus encore d'une véritable concorde entre la population d'une communauté et l'institution détentrice du pouvoir. C'est l'aspect le plus important chez Ambroise, celui qui conditionne toute la vie en société et qui doit donc être l'objectif principal du bon gouvernement. C'est pourquoi les divers exemples qu'il utilise ne se concentre pas seulement sur la forme que prend l'exercice du pouvoir, mais sur le lien entre les membres de la communauté et ce pouvoir. Dans ce cadre, l'analyse de l'\latin{Exameron} d'Ambroise a permis de mettre en lumière un aspect original de sa pensée : l'existence d'un idéal théorique de la \latin{res publica}. Ce modèle apparait clairement dans ses travaux comme le meilleur moyen d'atteindre la concorde. À travers l'exemple du vol ordonné des grues et du partage inné des tâches et pouvoirs, Ambroise valorise le modèle d'une société où le pouvoir n'est pas le privilège d'un seul, mais une charge répartie au service du bien commun. Selon l'évêque, cet agencement est celui de la nature voulue par Dieu, il est donc normal d'y trouver une forme optimale de l'entente dans une société.

\bigskip

Pourtant, il ne dévalue pas pour autant le modèle monarchique. Au-delà du régime impérial romain dont il fait partie, Ambroise utilise l'exemple des abeilles. Bien que vivant sous l'autorité d'un seul, elles parviennent à avancer dans une même direction, avce comme seul objecif la réussite de leur communauté. Et c'est là le point le plus important pour Ambroise : si le pouvoir est détenu par un monarque, alors celui-ci doit s'assurer de la confiance et du dévouement de ses sujets. Le problème de la monarchie survient lorsque le détenteur du pouvoir agit pour sa propre réussite, ce qui fait alors apparaître la tyrannie et la coercition. Derrière cette réflexion, Ambroise apporte deux aspects concrets : la nécessité pour le monarque d'agir pour le bien commun, et l'obligation pour la population de suivre, presque aveuglément, ce pouvoir afin d'aboutir à une parfaite concorde. Loin d'imposer un rejet de la monarchie, il est confronté à la légitimité du pouvoir impérial et décide de l'utiliser comme un outil pouvant permettre d'atteindre la \latin{concordia}, ainsi qu'une diffusion universelle de la foi chrétienne. Ambroise décide de briser la séparation intelectuelle entre monarchie et bien commun. Il place ainsi le souverain comme un garant de l'ordre social qui doit s'inspirer des vertus républicaines et chrétiennes pour consolider son dévouement envers l'Empire.

\bigskip

Toute cette réflexion ambrosienne est guidée par la pensée chrétienne de l'évêque. Il conçoit la possibilité du régime monarchique uniquement sous le prisme d'un encadrement chrétien plus large. Le monarque oeuvrant pour la concorde et le bien commun ne peut être que chrétien, ce qui est d'ailleurs légitimé par une vision chrétienne, par Dieu, de l'origine du pouvoir. Cette conception chrétienne du monde qui l'entoure se ressent évidemment dans ses textes et arguments. Il s'appuie en grande majorité sur des exemples scripturaires et sur des figures royales bibliques pour instaurer le modèle du bon monarque, et transorme les idées issues de penseurs païens en propos chrétiens.

\bigskip

La modélisation éthique, spirituelle et morale du prince, par des idées chrétiennes est ce qui a constituté le second axe majeur de notre étude. Puisque le système impérial est accepté comme une réalité temporel nécessaire, Ambroise cherche avant tout à le classer, le codifier, et le ranger derrière les exigences de la moralité chrétienne pour s'assurer que l'exercice du pouvoir ne bascule pas dans la tyrannie. Il instaure ainsi, dans une approche théorique au départ, des règles majeurs à suivre pour les empereurs romains. Il n'est plus question de se questionner sur le pouvoir en général, mais bien de voir ce qu'il est possible d'apporter de façon concrète sur les empereurs. C'est dans ce cadre qu'il revient le plus largement sur le concept de justice et ce qu'il représente chez lui. Il se détache d'une justice punitive par la force et la capacité de décision qu'offre l'\latin{imperium} pour imposer la vision d'une justice qui rejette la violence inutile et valorise la tempérance et la piété. Il est difficile de juger l'impact qu'a eu Ambroise sur les actions impériales : il n'a pas modelé de toute pièce le prince chrétien mais il a sans aucun doute poussé les empereurs à penser leurs actions selon une logique chrétienne. C'est notamment le cas du courage militaire et de la guerre. Tout en valorisant les figures non-violentes, il utilise son pragmatisme habituel pour accepter les conflits armés si ils servent la foi chrétienne. C'est donc dans une unique perpsective chrétienne qu'il conçoit le pouvoir politique et le gouvernement romain. En tant que représentant de la foi et de l'Église, il légitime son droit à modeler les empereurs selon sa propre volonté.

\bigskip

Ce travail, que l'on peut qualifier de pédagogique et de politique envers les empereurs se réalise à travers une utilisation de la Bible qui quitte le domaine spirituelle pour devenir un outil de réflexion sur le gouvernement. Principalement par la figure de David, il présente aux empereurs un «~miroir des princes~» pour guider les actions impériales. Son utilisation de David est l'un de ses plus grands coup de force théorique : il démontre que la grandeur du souverain ne réside pas dans son infaillibilité, mais dans sa capacité à reconnaître ses fautes. Le bon empereur est celui qui fait preuve de piété et qui est capable de faire pénitence, ce que Théodose finit par accepter. la rédemption de Théodose après le massacre de Thessalonique est certainement le meilleur symbole de l'impact d'Ambroise sur la politique romaine. Le résultat de la pénitence bénéficie aussi bien à Ambroise qu'à l'empereur. Par cette requête, l'évêque de Milan s'ancre définitivement comme une figure centrale de la cour impériale, tandis que Théodose consolide son autorité et son règne. Par la validation chrétienne, il s'assure le soutien d'une grande partie de la population et assoit ainsi son pouvoir malgré sa faute importante. D'une façon plus générale, sa volonté de modeler les empereurs chrétiens aboutit à de nouveaux axes de légitimation pour les souverains : l'exemplarité morale, le suivi de la foi dès le plus jeune âge et le respect des volontés divines.

\bigskip

Enfin, le troisième thème fondamental de notre étude concernait la mise en place de cette théorie dans le monde politique en s'attaquant aux différents domaines d'autorités. L'analyse des traités d'Ambroise et la connaissance de ses interventions de terrains démontre que sa pensée politique atteint son sommet lorsqu'elle s'articule avec l'action. Nous avons pu constater qu'à travers des évènements comme le conflit concernant le retour de l'autel de la Victoire dans le Sénat, les tensions autour de l'utilisation des basiliques de Milan ou encore les missions d'ambassades auprès de Maxime, Ambroise fait travailler l'autorité du rôle d'évêque. À travers ses travaux, et surtout ses lettres aux empereurs, il est possible de voir que l'évêque cherche à donner à l'Église chrétienne une \latin{auctoritas} qui permet à l'institution d'imposer son point de vu sur les questions religieuses, tout en conservant une forme d'autonomie vis-à-vis du pouvoir impérial. Puisque son questionnement théorique se fait à travers un cadre chrétien, que ce soit dans la priorité des actions, dans le suivi de la foi ou dans la restriction morale, il parrait assez logique que sa vision de la politique romaine se fasse avec une autorité chrétienne en son coeur. Il offre ainsi aux évêques, en se mettant lui-même en tête de file, la possibilité d'obtenir un rôle indispenable dans la vie temporel de l'Empire.

\bigskip

Mais ce qu'il faut surtout retenir de cette autorité épiscopale, est le fait qu'Ambroise use de sa réflexion théorique pour modeler son propre rôle auprès des empereursen engrageant plusieurs rôle : l'évêque officiel de la capitale impériale, l'aristrocrate politicien ou le conseiller politique et spirituel. L'arme majeur qu'il parvient à utiliser à ce sujet est celle de la \latin{parrhèsia}, ce franc-parler hérité de la philosophie antique, qu'Ambroise transforme en un devoir sacerdotal. C'est à partir de ça qu'il se permet de s'opposer aux décisions impériales, ou plus simplement de conseiller avec virulence les empereurs. Il est conscient que sa parole peu influencé de façon concrète la politique de l'Empire et il n'hésite pas à en jouer. La pensée politique d'Ambroise est donc intimement lié à son autorité et à celle de la foi chrétienne. En appuyant sur la distinction entre le domaine sacré et les affaires temporelles, il légitime sa présence et son activité sans pour autant restreindre le pouvoir du monarque. Ambroise a ainsi entamé le processus de développpement de l'autorité et de l'autonomie des figures religieuses que l'on constate ensuite avec les pouvoirs politiques et l'influence des papes en Occident. La volonté de restriction morale des détenteurs du pouvoir pour faciliter le bien commun n'est donc pas resté au stade d'utopie théorique mais commence avec Ambroise à devenir une réalité politique.

\bigskip

Au-delà de la réorganisation morale, administrative, intelectuelle ou encore juridique du pouvoir impérial, la véritable réussite d'Ambroise en tant que théologien et penseur politique réside dans la remise en question du statut de l'empereur au sein du christianisme. Sa réfléxion, même dans l'idéal théorique, ne s'attaque jamais frontalement au régime impérial et à la légitimité des empereurs. Ce qui compte vraiment pour l'évêque de Milan, c'est que le pouvoir romain soit soumis aux exigences chrétiennes, aussi bien dans la défense de la foi, dans la priorité donné aux affaires religieuses que dans le suivi d'une certaine morale et justice. Le spectre du christianisme doit résonner partout, y compris dans la nouvelle autorité des évêques. Ambroise est alors le meilleur représentant de la relation de la foi chrétienne avec les empereurs. Cette façon d'approcher la religion ne dessere par le pouvoir politique, au contraire elle lui offre une légitimité pour de nobreuses actions. Puisqu'AMbroise veut que les empereurs servent l'Église, il est logique que l'Église favories l'Empire. La lutte armée par exemple est souvent légitimée par des raisons de foi, tandis que l'essence même du pouvoir du souverain conforté par l'appui des évêques.

\bigskip

On est, à partir d'Ambroise, bien loin du modèle constantinien prôné par Eusèbe de Césarée qui voulait que l'emepreur soit ministre de Dieu et qu'il possède donc un rôle direct dans la diffusion du message de Dieu. Ambroise sépare distinctement l'autorité politique temporelle de l'autorité spirituelle des évêques qui, comme le Sénat, jouent un rôle d'influence auprès des décisions impériales. L'évêque de Milan est parvenu par ses traités et correspondances à encadrer le pouvoir des princes chrétiens sans jamais s'en prendre aux fondements de l'\latin{imperium}. Il ne cherche pas à s'opposer à la fonction impériale, mais à la modeler dans une nouvelle direction, sous l'autorité de la foi chrétienne. C'est une idée qui se résume particulièrement bien à travers une phrase que l'on retrouve dans son \latin{Epistola} 75A, «~Contre Auxence~». Cette lettre est la mise à l'écrit d'un sermon qu'il proclamme devant le peuple de Milan, il doit donc mettre en avant des arguments percutants pour montrer son soutien total envers la frange nicéenne de la population milanaise. Il insiste alors sur le fait que l'empereur est subordonné à l'Église chrétienne : «~\latin{Imperator enim intra ecclesiam non supra ecclesiam est}~»\autocite[\nopp 75A, 36]{ambroise_lettres}.

\bigskip

Il ne s'agit pas seulement d'un argument rhétorique, c'est l'expression de sa pensée politique et de sa vision de la relation de l'empereur avec l'Église. Pour autant, cette formule ne place pas l'Église et les évêques au-dessus du pouvoir impérial. Ambroise s'attache simplement à créer une distinction entre la fonction publique du souverain et son statut personnel de fidèle.

Il n'a pas fait des éveques les détenteurs d'un pouvoir d'action plus frt, mais il a donné une autotié insitutionellete et morale autonome capable de juger et donc de dicter certaines actions du souverian.

Et c'est à travers sa proximité avec le christianisme que l'empereur doit suivre les conseils des évêques et doit respecter les exigences que requièrent la foi. Ainsi, AMbroise de Milan est parvenu à «~christianisé~» le pouvoir romain tout autant qu'il a «~romanisé~» l'Église. Les deux institutuions doivent se respecter et s'entraider pour atteindre l'objectif de bien commun, qui est ce que recherche activement Ambroise.





Cette formule n'est pas une simple boutade rhétorique ; elle constitue une révolution théologico-politique. En déclarant que l'empereur est "dans" l'Église, Ambroise opère une séparation décisive entre la fonction publique du souverain et son statut personnel de fidèle. En tant que dirigeant de la \latin{res publica}, l'empereur détient un pouvoir politique absolu que l'évêque respecte scrupuleusement. Mais en tant que chrétien, l'homme revêtu de la pourpre est soumis aux mêmes exigences de Salut, à la même discipline pénitentielle et aux mêmes obligations morales que le plus humble de ses sujets. L'affaire de Thessalonique en est l'illustration la plus éclatante : lorsque Théodose accepte la pénitence publique, il ne s'incline pas en tant qu'empereur cédant son pouvoir politique, il s'incline en tant que pécheur implorant la miséricorde divine par l'intermédiaire de l'évêque.

\bigskip

C'est ici que se trouve le génie politique d'Ambroise de Milan. Il n'a pas cherché à transformer l'Empire en une théocratie où les évêques gouverneraient l'État. Il a au contraire "romanisé" l'Église en lui conférant une autorité institutionnelle et morale autonome, capable de juger les actes du souverain. Ce faisant, il a créé un espace sacré inviolable, indépendant du pouvoir temporel. Cette conception d'une Église gardienne de la morale, limitant l'arbitraire du pouvoir impérial sans chercher à le remplacer, pose les premiers jalons intellectuels et juridiques de ce qui deviendra la grande théologie politique du Moyen Âge. Ambroise a ainsi doté le christianisme latin de l'outil conceptuel le plus puissant de son histoire : une autorité spirituelle capable de tenir tête à l'absolutisme temporel.






(L'IDÉE MAÎTRESSE : Le couronnement de ton mémoire)
Cette dynamique nous amène à l'idée maîtresse qui couronne l'action ambrosienne : l'évêque de Milan a définitivement enterré le modèle politico-religieux constantinien et eusébien. Depuis Constantin, l'Empire concevait l'empereur comme «~l'évêque du dehors~», un représentant direct de Dieu sur terre, quasi chef de l'Église. Ambroise renverse totalement ce paradigme. Par sa théorie et ses actions, il instaure un nouveau rapport de force où, selon sa célèbre formule, «~\latin{Imperator intra Ecclesiam, non supra Ecclesiam est}~» (L'empereur est dans l'Église, et non au-dessus de l'Église). C'est ici la plus grande victoire d'Ambroise : il ne s'est pas contenté de christianiser le droit romain, il a "romanisé" l'Église en lui conférant une autorité institutionnelle capable de juger l'État. Il pose ainsi, plusieurs siècles à l'avance, les fondations intellectuelles et juridiques de ce qui deviendra la théologie politique du Moyen Âge, prémices de la lutte du Sacerdoce et de l'Empire.

(L'OUVERTURE : Vers Augustin et la fin de l'Empire)
Cependant, l'édifice intellectuel d'Ambroise comportait une faille inhérente à son propre patriotisme. Toute sa pensée politique reposait sur l'idée providentielle que le christianisme et l'Empire romain étaient intrinsèquement liés, l'un assurant la prospérité de l'autre. Ambroise est mort avant de voir cet idéal s'effondrer. C'est à son plus célèbre disciple intellectuel (qu'il a lui-même baptisé à Milan en 387), Augustin d'Hippone, qu'il reviendra de repenser ce modèle. Moins de quinze ans après la mort d'Ambroise, le sac de Rome par Alaric en 410 provoque une onde de choc qui ruine l'espoir d'une romanité chrétienne invulnérable. Face à cette catastrophe, Augustin devra rompre avec l'optimisme impérial de son maître. En rédigeant La Cité de Dieu, il séparera définitivement le destin de l'Église de celui de l'Empire romain, montrant que si les États terrestres sont mortels, seule la Cité céleste demeure. Si Ambroise a théorisé l'Église face au pouvoir de l'Empire, Augustin théorisera l'Église survivant à sa chute, achevant ainsi de détacher la pensée politique chrétienne des ruines de la romanité.
